\subsection{QC-011: full padded multiplicity through degree nine}

\paragraph{Statement QC-011.}
Let $E$ be a finite-dimensional complex vector space, let $n\geq1$,
and let $V_{n-1,n}\subseteq\operatorname{Sym}^nE$ be the closure of
the forms $LQ$, where $L\in E$ and $Q\in\operatorname{Sym}^{n-1}E$.
Then
\[
       I(V_{n-1,n})_d=0\qquad(0\leq d\leq n).
       \tag{QC-011.1}
\]
In particular, for the pair $(m,n)=(8,9)$ in the question,
\[
       q_\lambda=a_\lambda\qquad
       \text{for every four-row }\lambda\vdash9d,
       \quad 0\leq d\leq9.
       \tag{QC-011.2}
\]
There is no restriction on the tail shape in this degree range.
The assertion compares the complete multiplicity spaces.
It gives no determinant-side rank or vanishing assertion.

\paragraph{Literature status.}
The directional-weight-shift method is established, and no priority
claim is made for it.  In C.~Ikenmeyer, \emph{On McKay's propagation
theorem for the Foulkes conjecture}, arXiv:1509.04957v1, Section~3,
equation~(3.2), the right factor is the common-one-factor map under
Schur--Weyl duality; Section~5, Claim~5.1 and its proof, printed
pp.~6--7, prove its injectivity in the stated range $a<b$.
Here $a=d$ and $b=n$.  The direct proof below includes the boundary
$d=n$ and spells out the actual padded-ideal consequence.
The Fischer-inner-product lemma itself is also explicitly given in
OpenAI, \emph{Quadratic stabilization of the canonical Foulkes--Howe
map}, 25 September 2026, Section~3, Lemma~2
(source label \texttt{lem:shift}); Section~4, Step~4 uses the related
operation of removing a common power.  We use only this elementary
lemma, whose complete proof is checked below, not that work's
full Foulkes--Howe theorem.
Kadish--Landsberg, arXiv:1204.4693v1, Theorem~1.7 and Section~3,
identify the padded ideal with the kernel of the generalized
Foulkes--Howe map; this identification is compatible with the direct
restriction argument here.  Its separate sufficient branch corrected
in S-009 is not used.

\paragraph{An elementary directional-derivative lemma.}
Let $U=(U_1,\ldots,U_r)$ and $Y=(Y_1,\ldots,Y_r)$ be two sets
of polynomial variables, possibly together with further variables.
On polynomials homogeneous of degrees $p$ and $q$ in $U$ and $Y$,
respectively, the operator
\[
               A=\sum_{j=1}^rY_j\frac{\partial}{\partial U_j}
\]
is injective whenever $p>q$.

Here is a direct proof, including the case of polynomial coefficients
in the further variables.  Give the full complex polynomial ring the
positive definite Hermitian inner product for which distinct monomials
are orthogonal and the squared norm of $X^\alpha$ is $\alpha!$.
Multiplication by a variable is adjoint to differentiation in that
variable.  Thus the adjoint of $A$ is
\[
               B=\sum_{j=1}^rU_j\frac{\partial}{\partial Y_j}.
\]
Direct multiplication of differential operators gives
\[
        BA-AB=\sum_jU_j\partial_{U_j}-\sum_jY_j\partial_{Y_j}.
\]
For a polynomial $f$ of bidegree $(p,q)$ it follows that
\[
       \|Af\|^2=\|Bf\|^2+(p-q)\|f\|^2.
       \tag{QC-011.3}
\]
If $Af=0$ and $p>q$, positivity forces $f=0$.  No assumption
about symmetry in the $U$ variables or in the further variables is
needed.

\paragraph{Proof of the whole-space restriction statement.}
The case $d=0$ is immediate.  Let $1\leq d\leq n$, and suppose
a homogeneous polynomial $P$ of coefficient degree $d$ vanishes
on every $LQ$.  Polarize $P$ to a symmetric $d$-linear form
\[
       T:(\operatorname{Sym}^nE)^d\longrightarrow\mathbb C,
       \qquad P(F)=T(F,\ldots,F).
\]
For each fixed $L$, polarizing the identity $P(LQ)=0$ in $Q$
gives
\[
               T(LQ_1,\ldots,LQ_d)=0
               \quad\text{for all }Q_1,\ldots,Q_d.
               \tag{QC-011.4}
\]
Introduce independent vector variables $U_1,\ldots,U_d,Y$ in $E$
and define the scalar polynomial
\[
             f(U_1,\ldots,U_d)=T(U_1^n,\ldots,U_d^n).
\]
It has degree $n$ in each $U_i$ and degree zero in $Y$.  Set
\[
               A_i=\sum_{j=1}^{\dim E}Y_j
                                \frac{\partial}{\partial (U_i)_j}.
\]
These operators commute, and (QC-011.4), with
$Q_i=U_i^{n-1}$, gives
\[
       A_d\cdots A_1 f
          =n^d T(YU_1^{n-1},\ldots,YU_d^{n-1})=0.
          \tag{QC-011.5}
\]
Before the $i$th operator is applied, the degree in its own variable
$U_i$ remains $n$, while the degree in the shared variable $Y$
is $i-1$.  Since $i\leq d\leq n$, we have $n>i-1$.
The lemma therefore makes every step injective.  Applying these
injections successively to (QC-011.5) gives $f=0$.

Pure $n$th powers span $\operatorname{Sym}^nE$ in characteristic
zero: expand a general linear combination to the $n$th power and
compare its coefficients.  By multilinearity, $f=0$ consequently
implies $T=0$, and hence $P=0$.  This proves (QC-011.1).

For the permanent pair $(8,9)$, the verified quaternary permanent
dominance identifies its four-variable padded image with
$V_{8,9}$.  Applying (QC-011.1) with $E=\mathbb C^4$ shows
that restriction is injective on the entire degree-$d$ ambient
space for $d\leq9$.  In characteristic zero its decomposition into
irreducible $\mathrm{GL}(E)$-modules is semisimple, so every
highest-weight multiplicity is preserved.  The four-variable
reduction identifies these with the four-row multiplicities in
the question, proving (QC-011.2).

\paragraph{Consequences and scope.}
Every degree-six nonic candidate in the exact ambient sieve has
its padded multiplicity equal to its computed ambient multiplicity.
For example, the types $(14,14,14,12)$, $(15,15,12,12)$ and
$(27,9,9,9)$ have padded multiplicities $3,4,13$, respectively.
Those numerical ambient values are separate character computations;
the proof here establishes the entire-space equality independently
of them.  Deciding a multiplicity obstruction in these types now
requires only the determinant-side comparison.

This uses the restriction to forms with one linear factor.  It does
not assert injectivity of the complete Chow Foulkes--Howe map.
The operator proof is valid in arbitrary vector-space dimension,
whereas the identification with the particular permanent image used
in (QC-011.2) is the checked four-variable statement.

\paragraph{Geometric dictionary and verification.}
Equation (QC-011.1) says that the homogeneous coherent ideal of the
image of the multiplication morphism
$E\times\operatorname{Sym}^{n-1}E\to\operatorname{Sym}^nE$
has zero graded piece through degree $n$.  Equivalently its graded
restriction map has zero kernel in those degrees.  The passage to
Schur multiplicities takes place in the characteristic-zero
representation category, with the group-scheme and Schur-functor
dictionary used in S-007.  No identification with a normalization's
global-section ring is needed.
The author checked the whole proof, including the adjoint convention,
commutator sign, degree at every derivative step, polarization in
the common factor, and the final pure-power spanning argument.
